% ==============================================================================
% 11_appendix_member1.tex
% Phụ lục Phụ trách: Thành viên 1
% ==============================================================================
% Nội dung:
% - Prompt tạo QA bằng ChatGPT.
% - Dataset schema (JSON cấu trúc hoàn chỉnh).
% - Hướng dẫn gán nhãn và tiêu chí review của chuyên gia.
% - Ví dụ minh họa câu hỏi thuộc 3 mức độ: Easy, Medium, Hard.
% ==============================================================================

\section{Benchmark Construction and Annotation Guidelines (Member 1)}
\label{sec:appendix_member1}

\subsection{QA Generation Prompt for ChatGPT}
Figure~\ref{fig:prompt_qa_gen} provides the exact prompt instruction used by ChatGPT during the initial generation phase of the benchmark dataset.

\begin{figure}[h]
\centering
\small
\fbox{
\begin{minipage}{0.95\columnwidth}
\textbf{System Instruction for QA Generation:} \\
\texttt{You are a senior Angular architect. Given the following technical documentation page, formulate high-quality question-answering pairs that developers would realistically search for.} \\
\\
\textbf{Guidelines:}
\begin{enumerate}
    \item Formulate one direct conceptual question and one implementation/code-based question.
    \item Extract the exact verbatim evidence span(s) from the document that justify the answer.
    \item Label the scope as \texttt{single-section}, \texttt{cross-section}, or \texttt{cross-document}.
    \item Output valid JSON adhering strictly to the benchmark schema.
\end{enumerate}
\end{minipage}
}
\caption{System prompt used for LLM-assisted QA generation.}
\label{fig:prompt_qa_gen}
\end{figure}

\subsection{Benchmark QA Record Schema}
Every evaluation record in the benchmark is structured according to the JSON specification in Figure~\ref{fig:dataset_schema}.

\begin{figure}[h]
\centering
\footnotesize
\begin{minipage}{0.95\columnwidth}
\begin{verbatim}
{
  "query_id": "ANG-REC-089",
  "difficulty": "medium",
  "scope": "cross-section",
  "query_text": "How to bind input signal and",
  "             handle transform options?",
  "gold_evidence_spans": [
    {
      "doc_id": "signal-inputs.md",
      "header_path": "Signals > Inputs",
      "text": "input() declares an input..."
    }
  ],
  "gold_answer": "In Angular signals...",
  "requires_multihop": true
}
\end{verbatim}
\end{minipage}
\caption{JSON schema of an annotated benchmark record.}
\label{fig:dataset_schema}
\end{figure}

\subsection{Expert Annotation and Review Guidelines}
Reviewers evaluated candidate pairs across three mandatory quality criteria:
\begin{itemize}
    \item \textbf{Criterion 1 (Semantic Faithfulness):} The reference answer must be fully derivable from the designated evidence spans alone.
    \item \textbf{Criterion 2 (Syntactic Boundary Completeness):} Evidence spans containing code must encompass the complete enclosing logical block (e.g., full function declaration or class decorator).
    \item \textbf{Criterion 3 (Decoupled Ground Truth):} Ground truth boundaries must be recorded as character offsets and heading breadcrumbs, completely independent of chunk indices.
\end{itemize}

\subsection{Representative Query Examples across Difficulty Tiers}
\begin{itemize}
    \item \textbf{Easy (Single-Section):} \textit{``What is the syntax for declaring an output using the output() function in Angular 17+?''} \\
    $\rightarrow$ Grounded within a single subsection of \texttt{outputs.md}.
    
    \item \textbf{Medium (Cross-Section):} \textit{``How do you configure an HTTP interceptor using provideHttpClient and withInterceptors?''} \\
    $\rightarrow$ Requires synthesizing the setup configuration from the root configuration section and the interceptor definition subsection in \texttt{http.md}.
    
    \item \textbf{Hard (Cross-Document):} \textit{``How does deferrable views (@defer) interact with server-side rendering (SSR) hydration?''} \\
    $\rightarrow$ Spans distinct documentation files: \texttt{defer.md} (template syntax) and \texttt{ssr.md} (hydration lifecycle).
\end{itemize}
